\documentclass[12pt]{article}
\usepackage{amsmath,amssymb,amsfonts,amsthm}
\usepackage{hyperref}
\usepackage{geometry}
\usepackage{authblk}
\geometry{margin=1in}

\newtheorem{theorem}{Theorem}[section]
\newtheorem{lemma}[theorem]{Lemma}
\newtheorem{proposition}[theorem]{Proposition}
\newtheorem{corollary}[theorem]{Corollary}
\theoremstyle{definition}
\newtheorem{definition}[theorem]{Definition}
\newtheorem{assumption}[theorem]{Assumption}
\theoremstyle{remark}
\newtheorem{remark}[theorem]{Remark}

\newcommand{\Hthree}{\mathbb{H}^3}
\newcommand{\R}{\mathbb{R}}
\newcommand{\Z}{\mathbb{Z}}
\newcommand{\bdry}{\partial\mathbb{H}^3}

\title{Three Rigidity Properties of Logarithmic Weight Spectra\\
for Cocompact Lattice Actions on $\mathbb{H}^3$}
\date{}
\author[1]{Marvin L.\ Gentry}
\affil[1]{Independent Researcher\\
  \texttt{drmlgentry@protonmail.com}\\
  ORCID: 0009-0006-4550-2663}

\begin{document}
\maketitle

\begin{abstract}
Let $\Gamma$ be a torsion-free cocompact lattice in
$\mathrm{PSL}(2,\mathbb{C})$, acting on hyperbolic 3-space $\mathbb{H}^3$
by orientation-preserving isometries. Fix a basepoint $0\in\mathbb{H}^3$
and a positive constant $\kappa>0$, and define the logarithmic weight
$q(\gamma)=\kappa\,d(0,\gamma\cdot0)$ for $\gamma\in\Gamma$.
We prove three structural properties of the resulting weight spectrum:
\begin{enumerate}
\item[(i)] \textbf{Uniqueness.} Every boundary direction
      $\xi\in\partial\mathbb{H}^3$ has at most one minimal representative
      in $\Gamma$; since $\Gamma$ is torsion-free, this representative is
      unique whenever it exists.
\item[(ii)] \textbf{Discreteness.} The weight spectrum
      $\{q(\gamma):\gamma\in\Gamma\}$ has a positive minimum value
      $\Delta=\kappa\,\ell_0$, where $\ell_0$ is the systole of the action.
      In particular, for any $R>0$, only finitely many elements satisfy
      $q(\gamma)\le R$.
\item[(iii)] \textbf{Coordinate rigidity.} Given a finite collection of
      boundary directions with pairwise angular separation at least
      $2\varepsilon_0>0$, any deformation of the directions by less than
      $\varepsilon_0$ preserves both the minimal representative for each
      direction and its integer coordinates under the Milnor--\v{S}varc
      quasi-isometry.
\end{enumerate}
All three properties follow from proper discontinuity, cocompactness,
torsion-freeness, and the classification of isometries of $\mathbb{H}^3$.
No arithmetic or dynamical assumptions are required.
\end{abstract}

\section{Introduction}

Discrete logarithmic hierarchies arise naturally in geometric contexts
such as warped product geometries, renormalisation group flows, and
holographic bulk-boundary correspondences. When such a hierarchy is
generated by a cocompact group action on hyperbolic space, it is natural
to ask what structural constraints the geometry imposes on the weight
spectrum. This note provides three such constraints, each elementary
but useful as a foundation for applications.

The setting is a torsion-free cocompact lattice
$\Gamma\subset\mathrm{PSL}(2,\mathbb{C})$ and the associated weight
function $q(\gamma)=\kappa\,d(0,\gamma\cdot0)$. The weight spectrum
$\mathcal{W}=\{q(\gamma):\gamma\in\Gamma\}$ is the object of study.

Our main results establish:
\begin{itemize}
\item \emph{Uniqueness} (Theorem~\ref{thm:uniqueness}): each ideal
      boundary direction is represented by at most one element of minimal
      weight.
\item \emph{Discreteness} (Lemma~\ref{lem:gap} and
      Corollary~\ref{cor:discrete}): the spectrum has a positive minimum
      value equal to $\kappa$ times the systole; it accumulates nowhere.
\item \emph{Coordinate rigidity} (Theorem~\ref{thm:rigidity}): small
      deformations of the boundary data leave the discrete structure
      unchanged.
\end{itemize}

These results are purely geometric. The proofs use only the cocompactness
and torsion-freeness of $\Gamma$, together with standard facts about
isometries of $\mathbb{H}^3$~\cite{ratcliffe,maskit,bridsonhaefliger}.

\section{Geometric Setting}

Throughout, $\mathbb{H}^3$ denotes real hyperbolic 3-space with the
standard metric of constant curvature $-1$, and
$\partial\mathbb{H}^3\cong\mathbb{S}^2$ its sphere at infinity.

\begin{assumption}\label{ass:action}
$\Gamma$ is a torsion-free group acting properly discontinuously and
cocompactly on $\mathbb{H}^3$ by orientation-preserving isometries.
Fix a basepoint $0\in\mathbb{H}^3$.
\end{assumption}

\begin{definition}[Logarithmic weight]\label{def:weight}
For a positive constant $\kappa>0$, the \emph{logarithmic weight} of
$\gamma\in\Gamma$ is $q(\gamma)=\kappa\,d(0,\gamma\cdot0)$.
The \emph{weight spectrum} of $(\Gamma,\kappa)$ is
$\mathcal{W}(\Gamma,\kappa)=\{q(\gamma):\gamma\in\Gamma\}$.
\end{definition}

\begin{definition}[Isometry types]\label{def:types}
An isometry $\delta\in\mathrm{Isom}(\mathbb{H}^3)$ is:
\emph{elliptic} if it fixes a point of $\mathbb{H}^3$;
\emph{parabolic} if it fixes exactly one point of $\partial\mathbb{H}^3$
and no point of $\mathbb{H}^3$;
\emph{hyperbolic} (or loxodromic) if it fixes exactly two points
of $\partial\mathbb{H}^3$.
\end{definition}

\begin{lemma}[No parabolics in cocompact lattices]\label{lem:nopar}
A cocompact lattice $\Gamma\subset\mathrm{Isom}(\mathbb{H}^3)$ contains
no parabolic elements. Every nontrivial $\gamma\in\Gamma$ is hyperbolic.
\end{lemma}
\begin{proof}
See~\cite[Corollary~12.3.5]{ratcliffe}: the limit set of a cocompact
group is all of $\partial\mathbb{H}^3$, so the quotient
$\Gamma\backslash\mathbb{H}^3$ is compact.  A parabolic element fixes
a point on $\partial\mathbb{H}^3$, and that fixed point would be a cusp
point of the quotient, contradicting compactness.
\end{proof}

\section{Uniqueness of Minimal Representatives}
\label{sec:uniqueness}

\begin{definition}\label{def:convergeto}
We say $\gamma\in\Gamma$ \emph{converges to} $\xi\in\partial\mathbb{H}^3$
if the geodesic ray from $0$ in the direction of $\gamma\cdot0$ has ideal
endpoint $\xi$.
\end{definition}

\begin{definition}[Minimal representative]\label{def:minimal}
$\gamma\in\Gamma$ is a \emph{minimal representative} for $\xi$ if it
converges to $\xi$ and
$q(\gamma)=\inf\{q(\gamma'):\gamma'\in\Gamma,\,\gamma'\text{ converges to }\xi\}$.
\end{definition}

\begin{remark}
Existence of a minimal representative for each $\xi$ follows from
cocompactness: since the orbit $\Gamma\cdot0$ is dense in
$\mathbb{H}^3\cup\partial\mathbb{H}^3$ (the full compactification) and
the weight $q$ is proper (finitely many elements of weight $\le R$ for
each $R$), the infimum is attained.
\end{remark}

\begin{theorem}[Uniqueness]\label{thm:uniqueness}
For every $\xi\in\partial\mathbb{H}^3$, there is at most one minimal
representative in $\Gamma$.
\end{theorem}

\begin{proof}
Suppose $\gamma_1,\gamma_2\in\Gamma$ are two minimal representatives for
$\xi$. Set $\delta=\gamma_1^{-1}\gamma_2$. Since both $\gamma_i$ converge
to $\xi$, the element $\delta$ fixes $\xi\in\partial\mathbb{H}^3$.
By Definition~\ref{def:types}, an isometry fixing a boundary point is
either elliptic or parabolic. By Lemma~\ref{lem:nopar}, $\Gamma$ contains
no parabolics, so $\delta$ is elliptic. Since $\Gamma$ is torsion-free
(Assumption~\ref{ass:action}), the only elliptic element is the identity.
Hence $\delta=\mathrm{id}$ and $\gamma_1=\gamma_2$.
\end{proof}

\section{Discreteness of the Weight Spectrum}
\label{sec:gap}

\begin{definition}[Systole]
The \emph{systole} of the action is
$\ell_0=\min_{\gamma\neq\mathrm{id}}d(0,\gamma\cdot0)>0$.
\end{definition}

The systole is positive because the orbit $\Gamma\cdot0$ is discrete
(proper discontinuity) and $\Gamma$ has no fixed points in $\mathbb{H}^3$
(torsion-freeness).

\begin{lemma}[Positive minimum weight]\label{lem:gap}
The weight spectrum has positive minimum value
$\Delta:=\kappa\,\ell_0>0$.
That is, $q(\gamma)\ge\Delta$ for all $\gamma\neq\mathrm{id}$.
\end{lemma}

\begin{proof}
For $\gamma\neq\mathrm{id}$, we have $d(0,\gamma\cdot0)\ge\ell_0$ by
definition of the systole. Hence $q(\gamma)=\kappa\,d(0,\gamma\cdot0)
\ge\kappa\,\ell_0=\Delta$.
\end{proof}

\begin{remark}
The weight spectrum can have coincidences: distinct elements $\gamma,\gamma'$
may satisfy $q(\gamma)=q(\gamma')$ (for instance, $\gamma$ and any of its
conjugates, or $\gamma$ and $\gamma^{-1}$, have the same displacement from
$0$ if $0$ lies on the axis of $\gamma$). Lemma~\ref{lem:gap} asserts only
that the smallest nonzero weight is $\Delta$, not that all weights are
distinct.
\end{remark}

\begin{corollary}[Uniform discreteness]\label{cor:discrete}
For any $R>0$, the set $\{\gamma\in\Gamma:q(\gamma)\le R\}$ is finite.
The weight spectrum $\mathcal{W}(\Gamma,\kappa)$ has no accumulation points.
\end{corollary}

\begin{proof}
$\{q(\gamma)\le R\}=\{\gamma:d(0,\gamma\cdot0)\le R/\kappa\}$, which is
the intersection of $\Gamma$ with a closed ball of radius $R/\kappa$
in $\mathbb{H}^3$. This intersection is finite by proper discontinuity.
\end{proof}

\section{Coordinate Rigidity under Deformations}
\label{sec:rigidity}

Since $\Gamma$ acts cocompactly on $\mathbb{H}^3$, it is finitely generated.
Fix a finite generating set $\mathcal{S}$ and the associated word metric
$d_\mathcal{S}$ on $\Gamma$.

\begin{proposition}[Milnor--\v{S}varc quasi-isometry~\cite{bridsonhaefliger}]
\label{prop:msvarc}
The orbit map $\phi\colon(\Gamma,d_\mathcal{S})\to(\mathbb{H}^3,d)$,
$\gamma\mapsto\gamma\cdot0$, is a quasi-isometry: there exist constants
$C_1,C_2,C_3>0$ such that
\[
C_1\,d_\mathcal{S}(\gamma,\gamma')-C_3
\;\le\;
d(\gamma\cdot0,\gamma'\cdot0)
\;\le\;
C_2\,d_\mathcal{S}(\gamma,\gamma')+C_3
\]
for all $\gamma,\gamma'\in\Gamma$.
\end{proposition}

\begin{definition}[Integer coordinates]\label{def:intcoords}
For $\gamma\in\Gamma$ with reduced word $w=s_{i_1}^{\epsilon_1}\cdots
s_{i_k}^{\epsilon_k}$ ($\epsilon_j=\pm1$), the \emph{integer coordinate}
$a_i(\gamma)=\sum_{j:\,i_j=i}\epsilon_j$ records the net count of
generator $s_i$. The vector $(a_1,\ldots,a_r)\in\Z^r$ depends on the
chosen word but is uniquely determined by $\gamma$ for a fixed generating
set and normal form.
\end{definition}

Now fix a finite collection of boundary directions
$\xi_1,\ldots,\xi_m\in\partial\mathbb{H}^3$ with pairwise angular separations
\[
\varepsilon_0 := \tfrac{1}{2}\min_{i\neq j}\,d_{\mathbb{S}^2}(\xi_i,\xi_j) > 0.
\]

\begin{definition}[Bounded deformation]\label{def:deform}
A \emph{deformation of magnitude $\varepsilon$} is a collection of maps
$\xi_i\mapsto\xi_i'$ with $d_{\mathbb{S}^2}(\xi_i,\xi_i')\le\varepsilon$
for all $i$. If $\varepsilon<\varepsilon_0$, the perturbed directions
$\xi_1',\ldots,\xi_m'$ remain pairwise distinct, and each $\xi_i'$ lies
strictly closer to $\xi_i$ than to any other $\xi_j$ ($j\neq i$).
\end{definition}

\begin{theorem}[Coordinate rigidity]\label{thm:rigidity}
Let $\gamma_i$ be the (unique) minimal representative for $\xi_i$, for
each $i$. For any deformation with $\varepsilon<\varepsilon_0$, the minimal
representative for $\xi_i'$ is again $\gamma_i$, and its integer coordinates
are unchanged.
\end{theorem}

\begin{proof}
Fix $i$. For any $\gamma\in\Gamma$ converging to $\xi_j$ with $j\neq i$,
the ideal endpoint of the ray from $0$ through $\gamma\cdot0$ lies at
angular distance at least $2\varepsilon_0$ from $\xi_i$, hence at angular
distance at least $\varepsilon_0$ from $\xi_i'$.
In contrast, $\gamma_i$ converges to $\xi_i$, and its ideal endpoint lies
at distance $\le\varepsilon<\varepsilon_0$ from $\xi_i'$.
Hence $\gamma_i$ still converges to $\xi_i'$ (in the sense of
Definition~\ref{def:convergeto}) and no element converging to any
$\xi_j'$, $j\neq i$, can converge to $\xi_i'$.
Since the orbit is discrete and the weight is proper
(Corollary~\ref{cor:discrete}), the minimal weight for $\xi_i'$ is
attained at $\gamma_i$. By Theorem~\ref{thm:uniqueness}, $\gamma_i$ is
the unique minimal representative for $\xi_i'$, and its integer coordinates
are the same.
\end{proof}

\section{Conclusion}

We have established three rigidity properties of the logarithmic weight
spectrum of any torsion-free cocompact lattice acting on $\mathbb{H}^3$:
uniqueness of minimal representatives, existence of a positive minimum
weight, and stability of the discrete structure under bounded deformations
of the boundary data. The proofs are elementary and self-contained; they
require no special arithmetic structure and apply equally to arithmetic
and non-arithmetic lattices.

These properties provide a rigorous geometric foundation for applications
in which discrete logarithmic hierarchies arise from cocompact hyperbolic
group actions, such as the derivation of discrete weight spectra in
symmetric space constructions~\cite{Gentry2026masses}.

\bibliographystyle{plain}
\bibliography{gentry-rigidity}

\end{document}
